% =========================================================
% Chapter 9: Rotation of Rigid Bodies
% POSN 1 Physics Preparation Notes
% Language: Thai
% Compiler: XeLaTeX
% =========================================================

\chapter{การหมุนของวัตถุแข็งเกร็ง}

\begin{center}
    {\Large \textbf{Rotation of Rigid Bodies}}\\[4pt]
    {\normalsize เอกสารเตรียมสอบคัดเลือก สอวน. ฟิสิกส์ ค่าย 1}\\[6pt]
    {\normalsize จัดทำโดย \textbf{Tames Thanakrit Damduan}}\\[2pt]
    {\footnotesize \textcopyright\ 2026 Tames Thanakrit Damduan. All rights reserved.}\\[8pt]
\end{center}

\section*{เป้าหมายของบทเรียน}

หลังจากเรียนบทนี้ นักเรียนควรทำสิ่งต่อไปนี้ได้

\begin{enumerate}
    \item เข้าใจปริมาณเชิงมุม ได้แก่ มุม การกระจัดเชิงมุม ความเร็วเชิงมุม และความเร่งเชิงมุม
    \item เชื่อมความสัมพันธ์ระหว่างปริมาณเชิงเส้นและปริมาณเชิงมุมได้
    \item ใช้สมการการหมุนเมื่อความเร่งเชิงมุมคงตัวได้
    \item เข้าใจโมเมนต์ความเฉื่อยในฐานะความเฉื่อยต่อการหมุน
    \item ใช้พลังงานจลน์ของการหมุนและการกลิ้งโดยไม่ไถลได้
    \item เข้าใจความหมายของทอร์กและใช้สมการ $\sum \tau = I\alpha$ ได้
    \item ใช้เงื่อนไขสมดุลสถิต $\sum F_x=0$, $\sum F_y=0$, และ $\sum \tau=0$ ได้
    \item แก้โจทย์ล้อกลิ้งลงพื้นเอียงและโจทย์สมดุลแบบ สอวน. ได้
\end{enumerate}

\begin{tcolorbox}[title=แนวคิดสำคัญของบทนี้]
การหมุนคล้ายการเคลื่อนที่แนวตรงมาก เพียงเปลี่ยนจาก $x,v,a,m,F$ ไปเป็น $\theta,\omega,\alpha,I,\tau$ ถ้าจำคู่เทียบเหล่านี้ได้ บทนี้จะเป็นการนำกฎของนิวตันและพลังงานกลับมาใช้ในรูปแบบการหมุน
\end{tcolorbox}

\newpage{}

% =========================================================
\section{คณิตศาสตร์ก่อนเรียน: มุมในหน่วยเรเดียน}

\subsection{ทำไมต้องใช้เรเดียน}

ในการหมุน เรามักใช้มุมเป็นเรเดียน เพราะทำให้สูตรระหว่างระยะส่วนโค้งกับมุมเขียนได้ง่ายที่สุด

ถ้าวัตถุเคลื่อนที่บนวงกลมรัศมี $r$ และกวาดมุม $\theta$ ระยะส่วนโค้งคือ

\[
s=r\theta
\]

สูตรนี้ใช้ได้เมื่อ $\theta$ อยู่ในหน่วยเรเดียน

\begin{tcolorbox}[title=นิยามของเรเดียน]
มุม $1$ เรเดียน คือมุมที่ทำให้ความยาวส่วนโค้งเท่ากับรัศมีของวงกลม

\[
\theta=\frac{s}{r}
\]
\end{tcolorbox}

\subsection{การแปลงองศาเป็นเรเดียน}

หนึ่งรอบเต็มคือ

\[
360^\circ=2\pi\,\text{rad}
\]

ดังนั้น

\[
180^\circ=\pi\,\text{rad}
\]

\[
90^\circ=\frac{\pi}{2}\,\text{rad}
\]

\[
60^\circ=\frac{\pi}{3}\,\text{rad}
\]

\[
30^\circ=\frac{\pi}{6}\,\text{rad}
\]

\begin{examplebox}{ตัวอย่าง: ระยะส่วนโค้ง}
จุดบนขอบล้อรัศมี $0.50\,\si{m}$ หมุนไปเป็นมุม $60^\circ$ จงหาระยะส่วนโค้ง

แปลงมุมก่อน

\[
60^\circ=\frac{\pi}{3}
\]

ใช้

\[
s=r\theta
\]

แทนค่า

\[
s=(0.50)\left(\frac{\pi}{3}\right)
\]

\[
s=\frac{\pi}{6}\,\si{m}
\]
\end{examplebox}

\begin{exercisebox}{แบบฝึกหัด 9.1}
\begin{enumerate}
    \item จงแปลง $45^\circ$ เป็นเรเดียน
    \item ล้อรัศมี $0.20\,\si{m}$ หมุนไป $2\pi$ เรเดียน จุดบนขอบล้อเคลื่อนที่เป็นระยะเท่าไร
    \item ถ้าระยะส่วนโค้งเท่ากับ $3\,\si{m}$ และรัศมีเท่ากับ $0.5\,\si{m}$ มุมที่กวาดได้เป็นกี่เรเดียน
\end{enumerate}
\end{exercisebox}

% =========================================================
\section{ปริมาณเชิงมุม}

\subsection{การกระจัดเชิงมุม}

ในบทการเคลื่อนที่แนวตรง เราใช้ตำแหน่ง $x$ และการกระจัด $\Delta x$

ในการหมุน เราใช้ตำแหน่งเชิงมุม $\theta$ และการกระจัดเชิงมุม

\[
\Delta \theta=\theta_f-\theta_i
\]

\begin{tcolorbox}[title=การกระจัดเชิงมุม]
\[
\Delta \theta=\theta_f-\theta_i
\]
\end{tcolorbox}

\subsection{ความเร็วเชิงมุม}

ความเร็วเชิงมุมเฉลี่ยนิยามเป็น

\[
\omega_{\text{avg}}=\frac{\Delta \theta}{\Delta t}
\]

ถ้าต้องการค่าขณะหนึ่ง ใช้ลิมิต

\[
\omega=\frac{d\theta}{dt}
\]

\begin{tcolorbox}[title=ความเร็วเชิงมุม]
\[
\omega_{\text{avg}}=\frac{\Delta \theta}{\Delta t}
\]

\[
\omega=\frac{d\theta}{dt}
\]
\end{tcolorbox}

หน่วยของ $\omega$ คือ

\[
\si{rad/s}
\]

\subsection{ความเร่งเชิงมุม}

ความเร่งเชิงมุมคืออัตราการเปลี่ยนของความเร็วเชิงมุม

\[
\alpha_{\text{avg}}=\frac{\Delta \omega}{\Delta t}
\]

และค่าขณะหนึ่งคือ

\[
\alpha=\frac{d\omega}{dt}
\]

\begin{tcolorbox}[title=ความเร่งเชิงมุม]
\[
\alpha_{\text{avg}}=\frac{\Delta \omega}{\Delta t}
\]

\[
\alpha=\frac{d\omega}{dt}
\]
\end{tcolorbox}

หน่วยของ $\alpha$ คือ

\[
\si{rad/s^2}
\]

\subsection{เครื่องหมายของปริมาณเชิงมุม}

โดยทั่วไปกำหนดให้การหมุนทวนเข็มนาฬิกาเป็นบวก และการหมุนตามเข็มนาฬิกาเป็นลบ

\[
\text{ทวนเข็มนาฬิกา} = +
\]

\[
\text{ตามเข็มนาฬิกา} = -
\]

แต่ถ้าโจทย์กำหนดทิศบวกต่างจากนี้ ให้ใช้ตามโจทย์

\begin{examplebox}{ตัวอย่าง: ความเร็วเชิงมุมเฉลี่ย}
ล้อหมุนจาก $\theta_i=2\,\text{rad}$ ไปเป็น $\theta_f=14\,\text{rad}$ ในเวลา $3\,\si{s}$ จงหา $\omega_{\text{avg}}$

\[
\omega_{\text{avg}}=\frac{\Delta \theta}{\Delta t}
\]

\[
\omega_{\text{avg}}=\frac{14-2}{3}
\]

\[
\boxed{\omega_{\text{avg}}=4\,\si{rad/s}}
\]
\end{examplebox}

% =========================================================
\section{สมการการหมุนเมื่อความเร่งเชิงมุมคงตัว}

\subsection{เทียบกับการเคลื่อนที่แนวตรง}

สมการการหมุนมีรูปเหมือนสมการการเคลื่อนที่แนวตรง

\begin{center}
\begin{tabular}{|c|c|}
\hline
\textbf{การเคลื่อนที่แนวตรง} & \textbf{การหมุน} \\
\hline
$x$ & $\theta$ \\
\hline
$v$ & $\omega$ \\
\hline
$a$ & $\alpha$ \\
\hline
\end{tabular}
\end{center}

เมื่อความเร่งเชิงมุมคงตัว จะได้

\[
\omega_f=\omega_i+\alpha t
\]

\[
\Delta\theta=\omega_i t+\frac{1}{2}\alpha t^2
\]

\[
\omega_f^2=\omega_i^2+2\alpha\Delta\theta
\]

\[
\Delta\theta=\frac{\omega_i+\omega_f}{2}t
\]

\begin{tcolorbox}[title=สมการการหมุนเมื่อ $\alpha$ คงตัว]
\[
\omega_f=\omega_i+\alpha t
\]

\[
\Delta\theta=\omega_i t+\frac{1}{2}\alpha t^2
\]

\[
\omega_f^2=\omega_i^2+2\alpha\Delta\theta
\]

\[
\Delta\theta=\frac{\omega_i+\omega_f}{2}t
\]
\end{tcolorbox}

\begin{examplebox}{ตัวอย่าง: ล้อเริ่มหมุนจากหยุดนิ่ง}
ล้อเริ่มหมุนจากหยุดนิ่งด้วยความเร่งเชิงมุมคงตัว $4\,\si{rad/s^2}$ เป็นเวลา $3\,\si{s}$ จงหา $\omega_f$ และ $\Delta\theta$

เริ่มจาก

\[
\omega_i=0
\]

\[
\alpha=4\,\si{rad/s^2}
\]

\[
t=3\,\si{s}
\]

หา $\omega_f$

\[
\omega_f=\omega_i+\alpha t
\]

\[
\omega_f=0+(4)(3)
\]

\[
\omega_f=12\,\si{rad/s}
\]

หามุมที่หมุนได้

\[
\Delta\theta=\omega_i t+\frac{1}{2}\alpha t^2
\]

\[
\Delta\theta=0+\frac{1}{2}(4)(3^2)
\]

\[
\Delta\theta=18\,\text{rad}
\]
\end{examplebox}

\begin{exercisebox}{แบบฝึกหัด 9.2}
\begin{enumerate}
    \item วัตถุเริ่มหมุนจากหยุดนิ่งด้วย $\alpha=2\,\si{rad/s^2}$ เป็นเวลา $5\,\si{s}$ จงหา $\omega_f$
    \item ล้อหมุนด้วย $\omega_i=10\,\si{rad/s}$ และมี $\alpha=-2\,\si{rad/s^2}$ จงหาเวลาที่ล้อหยุด
    \item วัตถุหมุนด้วย $\omega_i=4\,\si{rad/s}$ และ $\alpha=3\,\si{rad/s^2}$ เป็นเวลา $2\,\si{s}$ จงหา $\Delta\theta$
\end{enumerate}
\end{exercisebox}

% =========================================================
\section{ความสัมพันธ์ระหว่างปริมาณเชิงเส้นและเชิงมุม}

\subsection{ตำแหน่งตามแนวสัมผัส}

ถ้าจุดหนึ่งอยู่ห่างจากแกนหมุนเป็นระยะ $r$ และหมุนไปมุม $\theta$ ระยะตามแนวโค้งคือ

\[
s=r\theta
\]

\subsection{ความเร็วแนวสัมผัส}

หาอนุพันธ์ของ

\[
s=r\theta
\]

เมื่อ $r$ คงตัว

\[
\frac{ds}{dt}=r\frac{d\theta}{dt}
\]

ดังนั้น

\[
v_t=r\omega
\]

\subsection{ความเร่งแนวสัมผัส}

หาอนุพันธ์ของ

\[
v_t=r\omega
\]

เมื่อ $r$ คงตัว

\[
a_t=r\alpha
\]

\subsection{ความเร่งสู่ศูนย์กลาง}

แม้ $\omega$ คงตัว ทิศของความเร็วก็เปลี่ยนตลอดเวลา จึงมีความเร่งสู่ศูนย์กลาง

\[
a_c=\frac{v_t^2}{r}
\]

เพราะ

\[
v_t=r\omega
\]

จึงได้

\[
a_c=\frac{(r\omega)^2}{r}
\]

\[
a_c=r\omega^2
\]

\begin{tcolorbox}[title=สรุปความสัมพันธ์เชิงเส้นและเชิงมุม]
\[
s=r\theta
\]

\[
v_t=r\omega
\]

\[
a_t=r\alpha
\]

\[
a_c=\frac{v_t^2}{r}=r\omega^2
\]
\end{tcolorbox}

\begin{examplebox}{ตัวอย่าง: จุดบนขอบล้อ}
ล้อรัศมี $0.40\,\si{m}$ หมุนด้วย $\omega=5\,\si{rad/s}$ และมี $\alpha=2\,\si{rad/s^2}$ จงหา $v_t$ และ $a_t$ ของจุดบนขอบล้อ

\[
v_t=r\omega
\]

\[
v_t=(0.40)(5)
\]

\[
v_t=2.0\,\si{m/s}
\]

ต่อไป

\[
a_t=r\alpha
\]

\[
a_t=(0.40)(2)
\]

\[
a_t=0.80\,\si{m/s^2}
\]
\end{examplebox}

% =========================================================
\section{โมเมนต์ความเฉื่อย}

\subsection{ความหมายของโมเมนต์ความเฉื่อย}

ในการเคลื่อนที่แนวตรง มวล $m$ บอกว่าวัตถุเปลี่ยนความเร็วได้ยากแค่ไหน

ในการหมุน ปริมาณที่บอกว่าวัตถุเปลี่ยนการหมุนได้ยากแค่ไหนคือโมเมนต์ความเฉื่อย

\[
I
\]

สำหรับอนุภาคมวล $m$ ที่อยู่ห่างแกนหมุนเป็นระยะ $r$

\[
I=mr^2
\]

สำหรับระบบหลายอนุภาค

\[
I=\sum_i m_i r_i^2
\]

\begin{tcolorbox}[title=โมเมนต์ความเฉื่อย]
\[
I=\sum_i m_i r_i^2
\]

โมเมนต์ความเฉื่อยขึ้นกับทั้งมวลและการกระจายมวลรอบแกนหมุน
\end{tcolorbox}

\subsection{ทำไมวัตถุรูปร่างต่างกันมี $I$ ต่างกัน}

มวลที่อยู่ไกลแกนหมุนทำให้ $I$ มาก เพราะมี $r^2$ คูณอยู่

ดังนั้นวัตถุสองชิ้นที่มีมวลเท่ากัน อาจมีโมเมนต์ความเฉื่อยต่างกันได้ ถ้ามวลกระจายต่างกัน

\begin{tcolorbox}[title=ข้อควรจำ]
มวลอยู่ใกล้แกนหมุนมาก หมุนง่าย

มวลอยู่ไกลแกนหมุนมาก หมุนยาก
\end{tcolorbox}

\subsection{โมเมนต์ความเฉื่อยที่ควรรู้}

ในข้อสอบ สอวน. มักให้ค่า $I$ มาในโจทย์ แต่ควรรู้ค่าพื้นฐานบางตัว

\begin{center}
\begin{tabular}{|p{0.44\linewidth}|p{0.36\linewidth}|}
\hline
\textbf{วัตถุและแกนหมุน} & \textbf{โมเมนต์ความเฉื่อย} \\
\hline
วงแหวนบาง รอบแกนผ่านศูนย์กลางตั้งฉากกับระนาบวงแหวน & $I=MR^2$ \\
\hline
จานตัน รอบแกนผ่านศูนย์กลางตั้งฉากกับจาน & $I=\frac{1}{2}MR^2$ \\
\hline
ทรงกลมตัน รอบแกนผ่านศูนย์กลาง & $I=\frac{2}{5}MR^2$ \\
\hline
แท่งยาว $L$ รอบแกนผ่านจุดกึ่งกลางและตั้งฉากกับแท่ง & $I=\frac{1}{12}ML^2$ \\
\hline
แท่งยาว $L$ รอบแกนผ่านปลายแท่งและตั้งฉากกับแท่ง & $I=\frac{1}{3}ML^2$ \\
\hline
\end{tabular}
\end{center}

\begin{examplebox}{ตัวอย่าง: เปรียบเทียบความยากในการหมุน}
จานตันและวงแหวนบางมีมวล $M$ และรัศมี $R$ เท่ากัน หมุนรอบแกนผ่านศูนย์กลางเดียวกัน วัตถุใดหมุนยากกว่า

จานตันมี

\[
I_{\text{disk}}=\frac{1}{2}MR^2
\]

วงแหวนมี

\[
I_{\text{ring}}=MR^2
\]

เพราะ

\[
I_{\text{ring}}>I_{\text{disk}}
\]

ดังนั้นวงแหวนหมุนยากกว่า
\end{examplebox}

\begin{exercisebox}{แบบฝึกหัด 9.3}
\begin{enumerate}
    \item อนุภาคมวล $2\,\si{kg}$ อยู่ห่างแกนหมุน $0.5\,\si{m}$ จงหา $I$
    \item อนุภาคสองตัวมวล $m$ เท่ากันอยู่ห่างแกน $R$ และ $2R$ จงหา $I_{\text{รวม}}$
    \item จานตันและทรงกลมตันมีมวลและรัศมีเท่ากัน อันใดมีโมเมนต์ความเฉื่อยน้อยกว่า
\end{enumerate}
\end{exercisebox}

% =========================================================
\section{พลังงานจลน์ของการหมุน}

\subsection{เริ่มจากพลังงานจลน์ของอนุภาค}

อนุภาคหนึ่งที่อยู่ห่างแกนหมุนเป็นระยะ $r_i$ มีความเร็ว

\[
v_i=r_i\omega
\]

พลังงานจลน์ของอนุภาคนั้นคือ

\[
K_i=\frac{1}{2}m_i v_i^2
\]

แทน $v_i=r_i\omega$

\[
K_i=\frac{1}{2}m_i(r_i\omega)^2
\]

\[
K_i=\frac{1}{2}m_i r_i^2\omega^2
\]

สำหรับวัตถุแข็งเกร็งทั้งหมด

\[
K_{\text{rot}}=\sum_i \frac{1}{2}m_i r_i^2\omega^2
\]

ดึง $\frac{1}{2}\omega^2$ ออก

\[
K_{\text{rot}}=\frac{1}{2}\left(\sum_i m_i r_i^2\right)\omega^2
\]

แต่

\[
I=\sum_i m_i r_i^2
\]

ดังนั้น

\[
K_{\text{rot}}=\frac{1}{2}I\omega^2
\]

\begin{tcolorbox}[title=พลังงานจลน์ของการหมุน]
\[
K_{\text{rot}}=\frac{1}{2}I\omega^2
\]
\end{tcolorbox}

\subsection{การกลิ้งโดยไม่ไถล}

ถ้าวัตถุกลิ้งโดยไม่ไถล จะมีความสัมพันธ์

\[
v_{\text{cm}}=\omega R
\]

และ

\[
a_{\text{cm}}=\alpha R
\]

พลังงานจลน์รวมของวัตถุกลิ้งคือผลรวมของพลังงานจลน์ของศูนย์กลางมวลและพลังงานจลน์ของการหมุนรอบศูนย์กลางมวล

\[
K=\frac{1}{2}Mv_{\text{cm}}^2+\frac{1}{2}I_{\text{cm}}\omega^2
\]

\begin{tcolorbox}[title=การกลิ้งโดยไม่ไถล]
\[
v_{\text{cm}}=\omega R
\]

\[
a_{\text{cm}}=\alpha R
\]

\[
K=\frac{1}{2}Mv_{\text{cm}}^2+\frac{1}{2}I_{\text{cm}}\omega^2
\]
\end{tcolorbox}

\begin{examplebox}{ตัวอย่าง: ล้อกลิ้งด้วยอัตราเร็ว $v$}
ล้อมีมวล $M$ รัศมี $R$ และโมเมนต์ความเฉื่อย $I_{\text{cm}}=\frac{1}{2}MR^2$ กลิ้งโดยไม่ไถลด้วยอัตราเร็วศูนย์กลางมวล $v$ จงหาพลังงานจลน์รวม

เริ่มจาก

\[
K=\frac{1}{2}Mv^2+\frac{1}{2}I_{\text{cm}}\omega^2
\]

เพราะกลิ้งโดยไม่ไถล

\[
\omega=\frac{v}{R}
\]

แทนค่า

\[
K=\frac{1}{2}Mv^2+\frac{1}{2}\left(\frac{1}{2}MR^2\right)\left(\frac{v}{R}\right)^2
\]

\[
K=\frac{1}{2}Mv^2+\frac{1}{4}Mv^2
\]

\[
\boxed{K=\frac{3}{4}Mv^2}
\]
\end{examplebox}

% =========================================================
\section{ทอร์กและกฎการหมุน}

\subsection{ทอร์กคืออะไร}

แรงที่ทำให้วัตถุหมุนได้ไม่ได้ขึ้นกับขนาดของแรงอย่างเดียว แต่ขึ้นกับตำแหน่งที่แรงกระทำด้วย

ถ้าออกแรงใกล้บานพับ ประตูหมุนยาก

ถ้าออกแรงไกลบานพับ ประตูหมุนง่าย

ปริมาณที่บอกความสามารถของแรงในการทำให้วัตถุหมุน เรียกว่า ทอร์ก หรือ โมเมนต์ของแรง

\[
\tau=rF\sin\phi
\]

โดยที่

\[
r=\text{ระยะจากแกนหมุนถึงจุดที่แรงกระทำ}
\]

\[
F=\text{ขนาดของแรง}
\]

\[
\phi=\text{มุมระหว่าง } \vec{r} \text{ และ } \vec{F}
\]

\begin{tcolorbox}[title=ทอร์ก]
\[
\tau=rF\sin\phi
\]

หรือ

\[
\tau=F d
\]

โดย $d$ คือระยะแขนโมเมนต์ ซึ่งเป็นระยะตั้งฉากจากแกนหมุนถึงแนวแรง
\end{tcolorbox}

หน่วยของทอร์กคือ

\[
\si{N.m}
\]

\subsection{เครื่องหมายของทอร์ก}

กำหนดให้ทอร์กที่ทำให้วัตถุหมุนทวนเข็มนาฬิกาเป็นบวก และตามเข็มนาฬิกาเป็นลบ

\[
\tau_{\text{CCW}}>0
\]

\[
\tau_{\text{CW}}<0
\]

แต่ถ้าโจทย์กำหนดทิศบวกเอง ให้ใช้ตามโจทย์

\subsection{กฎข้อที่สองสำหรับการหมุน}

ในการเคลื่อนที่แนวตรง

\[
\sum F=ma
\]

ในการหมุนรอบแกนคงที่

\[
\sum \tau=I\alpha
\]

\begin{tcolorbox}[title=กฎการหมุน]
\[
\sum \tau=I\alpha
\]

ทอร์กลัพธ์ทำให้เกิดความเร่งเชิงมุม
\end{tcolorbox}

\begin{examplebox}{ตัวอย่าง: ทอร์กจากแรงตั้งฉาก}
แท่งยาวเบาหมุนได้รอบปลายด้านหนึ่ง มีแรง $20\,\si{N}$ กระทำตั้งฉากที่ระยะ $0.50\,\si{m}$ จากแกนหมุน จงหาขนาดทอร์ก

\[
\tau=rF\sin90^\circ
\]

\[
\tau=(0.50)(20)(1)
\]

\[
\boxed{\tau=10\,\si{N.m}}
\]
\end{examplebox}

\begin{examplebox}{ตัวอย่าง: ความเร่งเชิงมุมจากทอร์ก}
ล้อมีโมเมนต์ความเฉื่อย $I=2.0\,\si{kg.m^2}$ ถูกกระทำด้วยทอร์กลัพธ์ $6.0\,\si{N.m}$ จงหาความเร่งเชิงมุม

ใช้

\[
\sum\tau=I\alpha
\]

ดังนั้น

\[
\alpha=\frac{\sum\tau}{I}
\]

\[
\alpha=\frac{6.0}{2.0}
\]

\[
\boxed{\alpha=3.0\,\si{rad/s^2}}
\]
\end{examplebox}

\begin{exercisebox}{แบบฝึกหัด 9.4}
\begin{enumerate}
    \item แรง $10\,\si{N}$ กระทำตั้งฉากที่ระยะ $0.30\,\si{m}$ จากแกนหมุน จงหาทอร์ก
    \item แรง $20\,\si{N}$ กระทำที่ระยะ $0.50\,\si{m}$ โดยทำมุม $30^\circ$ กับแนวแขนโมเมนต์ จงหาทอร์ก
    \item วัตถุมี $I=4\,\si{kg.m^2}$ และมีทอร์กลัพธ์ $12\,\si{N.m}$ จงหา $\alpha$
    \item ถ้าออกแรงเท่าเดิมแต่เพิ่มแขนโมเมนต์เป็นสองเท่า ทอร์กเปลี่ยนอย่างไร
\end{enumerate}
\end{exercisebox}

% =========================================================
\section{ล้อกลิ้งลงพื้นเอียง}

\subsection{แรงที่เกี่ยวข้อง}

พิจารณาล้อรัศมี $R$ มวล $M$ โมเมนต์ความเฉื่อยรอบศูนย์กลางมวล $I$ กลิ้งโดยไม่ไถลลงพื้นเอียงมุม $\theta$

แรงที่กระทำมี

\[
Mg
\]

\[
N
\]

และแรงเสียดทานสถิต

\[
f
\]

แรงเสียดทานสถิตทำหน้าที่สร้างทอร์กให้ล้อหมุน

\begin{tcolorbox}[title=ข้อควรจำ]
ในการกลิ้งโดยไม่ไถลบนพื้นเอียง แรงเสียดทานเป็นแรงเสียดทานสถิต และไม่จำเป็นต้องทำงานสุทธิ เพราะจุดสัมผัสไม่ไถลเทียบกับพื้น
\end{tcolorbox}

\subsection{สมการตามแนวพื้นเอียง}

เลือกทิศลงตามพื้นเอียงเป็นบวก

แรงตามแนวพื้นเอียงคือ

\[
Mg\sin\theta
\]

ลงตามพื้น และแรงเสียดทาน $f$ ขึ้นตามพื้น

ดังนั้น

\[
Mg\sin\theta-f=Ma
\]

\subsection{สมการทอร์ก}

เลือกทิศหมุนลงพื้นเอียงเป็นบวก

ทอร์กที่ทำให้ล้อหมุนเกิดจากแรงเสียดทาน

\[
fR=I\alpha
\]

และเงื่อนไขกลิ้งโดยไม่ไถลคือ

\[
a=\alpha R
\]

ดังนั้น

\[
\alpha=\frac{a}{R}
\]

แทนในสมการทอร์ก

\[
fR=I\frac{a}{R}
\]

\[
f=\frac{I}{R^2}a
\]

แทนลงในสมการแรง

\[
Mg\sin\theta-\frac{I}{R^2}a=Ma
\]

\[
Mg\sin\theta=a\left(M+\frac{I}{R^2}\right)
\]

ดังนั้น

\[
a=\frac{Mg\sin\theta}{M+\frac{I}{R^2}}
\]

หรือ

\[
a=\frac{g\sin\theta}{1+\frac{I}{MR^2}}
\]

และแรงเสียดทานคือ

\[
f=\frac{I}{R^2}a
\]

จึงได้

\[
f=
\frac{I}{R^2}
\left(
\frac{Mg\sin\theta}{M+\frac{I}{R^2}}
\right)
\]

จัดรูป

\[
f=\frac{I}{MR^2+I}Mg\sin\theta
\]

\begin{tcolorbox}[title=ล้อกลิ้งลงพื้นเอียงโดยไม่ไถล]
\[
a=\frac{g\sin\theta}{1+\frac{I}{MR^2}}
\]

\[
f=\frac{I}{MR^2+I}Mg\sin\theta
\]

แรงเสียดทานมีทิศขึ้นตามพื้นเอียง
\end{tcolorbox}

\begin{examplebox}{ตัวอย่าง: จานตันกลิ้งลงพื้นเอียง}
จานตันมี $I=\frac{1}{2}MR^2$ กลิ้งโดยไม่ไถลลงพื้นเอียงมุม $\theta$ จงหา $a$

ใช้

\[
a=\frac{g\sin\theta}{1+\frac{I}{MR^2}}
\]

แทน

\[
I=\frac{1}{2}MR^2
\]

จะได้

\[
a=\frac{g\sin\theta}{1+\frac{1}{2}}
\]

\[
a=\frac{g\sin\theta}{\frac{3}{2}}
\]

\[
\boxed{a=\frac{2}{3}g\sin\theta}
\]
\end{examplebox}

% =========================================================
\section{สมดุลสถิต}

\subsection{เงื่อนไขสมดุล}

วัตถุแข็งเกร็งจะอยู่ในสมดุลสถิตเมื่อไม่มีการเลื่อนที่และไม่มีการหมุน

ดังนั้นต้องมีทั้ง

\[
\sum F_x=0
\]

\[
\sum F_y=0
\]

และ

\[
\sum \tau=0
\]

\begin{tcolorbox}[title=เงื่อนไขสมดุลสถิต]
\[
\sum F_x=0
\]

\[
\sum F_y=0
\]

\[
\sum \tau=0
\]
\end{tcolorbox}

\subsection{เลือกจุดหมุนอย่างไร}

ในการตั้งสมการทอร์ก เราสามารถเลือกจุดหมุนใดก็ได้ ถ้าวัตถุอยู่ในสมดุลจริง ทอร์กรวมรอบทุกจุดต้องเป็นศูนย์

แต่เลือกจุดหมุนให้ดีจะทำให้สมการง่าย เช่น เลือกจุดที่มีแรงไม่รู้ค่าผ่านจุดนั้น เพื่อให้แรงนั้นไม่สร้างทอร์ก

\begin{tcolorbox}[title=เทคนิคเลือกจุดหมุน]
ถ้ามีแรงที่ไม่รู้ค่าหลายแรง ให้เลือกจุดหมุนที่แนวแรงไม่รู้ค่าผ่านจุดนั้น เพราะทอร์กของแรงนั้นจะเป็นศูนย์
\end{tcolorbox}

\subsection{ศูนย์กลางมวลและน้ำหนัก}

น้ำหนักของวัตถุแข็งเกร็งสามารถถือว่ากระทำที่ศูนย์กลางมวล

สำหรับแท่งสม่ำเสมอยาว $L$ ศูนย์กลางมวลอยู่ที่กึ่งกลางแท่ง

\[
x_{\text{cm}}=\frac{L}{2}
\]

\begin{examplebox}{ตัวอย่าง: แท่งวางบนที่รองรับ}
แท่งสม่ำเสมอยาว $L$ หนัก $W$ วางบนจุดรองรับสองจุดที่ปลายทั้งสอง จงหาแรงปฏิกิริยาที่แต่ละจุดรองรับ

ให้แรงที่ปลายซ้ายเป็น $N_1$ และปลายขวาเป็น $N_2$

สมดุลแนวดิ่ง

\[
N_1+N_2-W=0
\]

เลือกจุดหมุนที่ปลายซ้าย

\[
\sum\tau=0
\]

แรง $N_1$ ไม่สร้างทอร์ก เพราะผ่านจุดหมุน

\[
N_2L-W\left(\frac{L}{2}\right)=0
\]

ดังนั้น

\[
N_2=\frac{W}{2}
\]

จาก

\[
N_1+N_2=W
\]

ได้

\[
N_1=\frac{W}{2}
\]

ดังนั้น

\[
\boxed{N_1=N_2=\frac{W}{2}}
\]
\end{examplebox}

% =========================================================
\section{ชุดโจทย์รวมท้ายบท: คละระดับสำหรับ สอวน. ค่าย 1}

\subsection*{ระดับพื้นฐาน}

\begin{enumerate}
    \item จงแปลง $120^\circ$ เป็นเรเดียน
    \item ล้อรัศมี $0.5\,\si{m}$ หมุนด้วย $\omega=4\,\si{rad/s}$ จงหาอัตราเร็วของจุดบนขอบล้อ
    \item วัตถุหมุนจากหยุดนิ่งด้วย $\alpha=3\,\si{rad/s^2}$ เป็นเวลา $4\,\si{s}$ จงหา $\omega_f$
    \item อนุภาคมวล $2\,\si{kg}$ อยู่ห่างแกนหมุน $3\,\si{m}$ จงหาโมเมนต์ความเฉื่อย
\end{enumerate}

\subsection*{ระดับกลาง}

\begin{enumerate}
    \setcounter{enumi}{4}
    \item จานตันมวล $M$ รัศมี $R$ หมุนด้วย $\omega$ จงหาพลังงานจลน์ของการหมุน
    \item แรง $F$ กระทำตั้งฉากที่ปลายแท่งยาว $L$ ซึ่งหมุนได้รอบปลายอีกด้าน จงหาทอร์ก
    \item วัตถุมีโมเมนต์ความเฉื่อย $I$ ถูกกระทำด้วยทอร์กคงตัว $\tau$ จากหยุดนิ่งเป็นเวลา $t$ จงหา $\omega$
    \item วงแหวนบางกลิ้งโดยไม่ไถลด้วยความเร็วศูนย์กลางมวล $v$ จงหาพลังงานจลน์รวม
\end{enumerate}

\subsection*{ระดับยาก}

\begin{enumerate}
    \setcounter{enumi}{8}
    \item ล้อรัศมี $R$ มวล $M$ มีโมเมนต์ความเฉื่อย $I$ กลิ้งโดยไม่ไถลลงพื้นเอียงมุม $\theta$ จงหาความเร่งของศูนย์กลางมวล
    \item จากข้อก่อนหน้า จงหาแรงเสียดทานที่พื้นเอียงกระทำต่อล้อ
    \item แท่งสม่ำเสมอยาว $L$ หนัก $W$ ถูกยึดด้วยบานพับที่ปลายซ้าย และมีเชือกดึงปลายขวาขึ้นในแนวดิ่ง จงหาแรงตึงเชือก
    \item ประตูหนัก $W$ กว้าง $b$ มีบานพับบนและล่างห่างกัน $h$ ถ้าแรงแนวระดับที่บานพับบนมีขนาด $H$ จงหา $H$ จากสมดุลทอร์ก
\end{enumerate}

% =========================================================
\section{เฉลยย่อท้ายบท}

\begin{enumerate}
    \item
    \[
    120^\circ=\frac{2\pi}{3}\,\text{rad}
    \]

    \item
    \[
    v_t=r\omega=(0.5)(4)=2\,\si{m/s}
    \]

    \item
    \[
    \omega_f=\omega_i+\alpha t=0+(3)(4)=12\,\si{rad/s}
    \]

    \item
    \[
    I=mr^2=(2)(3^2)=18\,\si{kg.m^2}
    \]

    \item
    สำหรับจานตัน
    \[
    I=\frac{1}{2}MR^2
    \]
    ดังนั้น
    \[
    K_{\text{rot}}=\frac{1}{2}I\omega^2
    =
    \frac{1}{2}\left(\frac{1}{2}MR^2\right)\omega^2
    \]
    \[
    K_{\text{rot}}=\frac{1}{4}MR^2\omega^2
    \]

    \item
    \[
    \tau=FL
    \]

    \item
    จาก
    \[
    \tau=I\alpha
    \]
    ได้
    \[
    \alpha=\frac{\tau}{I}
    \]
    จากหยุดนิ่ง
    \[
    \omega=\alpha t
    \]
    ดังนั้น
    \[
    \omega=\frac{\tau t}{I}
    \]

    \item
    วงแหวนบางมี
    \[
    I=MR^2
    \]
    และ
    \[
    \omega=\frac{v}{R}
    \]
    ดังนั้น
    \[
    K=\frac{1}{2}Mv^2+\frac{1}{2}MR^2\left(\frac{v}{R}\right)^2
    \]
    \[
    K=Mv^2
    \]

    \item
    ใช้ผลจากล้อกลิ้งลงพื้นเอียง
    \[
    a=\frac{g\sin\theta}{1+\frac{I}{MR^2}}
    \]

    \item
    \[
    f=\frac{I}{R^2}a
    \]
    แทน $a$
    \[
    f=
    \frac{I}{R^2}
    \left(
    \frac{Mg\sin\theta}{M+\frac{I}{R^2}}
    \right)
    \]
    \[
    f=\frac{I}{MR^2+I}Mg\sin\theta
    \]

    \item
    เลือกจุดหมุนที่บานพับ
    \[
    TL-W\left(\frac{L}{2}\right)=0
    \]
    \[
    T=\frac{W}{2}
    \]

    \item
    เลือกจุดหมุนที่บานพับล่าง
    น้ำหนักประตูกระทำที่ระยะ $b/2$ จากแนวบานพับ และแรงแนวระดับที่บานพับบนมีแขนโมเมนต์ $h$
    \[
    Hh=W\left(\frac{b}{2}\right)
    \]
    ดังนั้น
    \[
    H=\frac{Wb}{2h}
    \]
\end{enumerate}

% =========================================================
\section{สรุปสูตรสำคัญ}

\begin{tcolorbox}[title=สรุปบท: การหมุนของวัตถุแข็งเกร็ง]
มุมและปริมาณเชิงมุม

\[
s=r\theta
\]

\[
\omega=\frac{d\theta}{dt}
\]

\[
\alpha=\frac{d\omega}{dt}
\]

เมื่อ $\alpha$ คงตัว

\[
\omega_f=\omega_i+\alpha t
\]

\[
\Delta\theta=\omega_i t+\frac{1}{2}\alpha t^2
\]

\[
\omega_f^2=\omega_i^2+2\alpha\Delta\theta
\]

ความสัมพันธ์เชิงเส้นและเชิงมุม

\[
v_t=r\omega
\]

\[
a_t=r\alpha
\]

\[
a_c=r\omega^2=\frac{v_t^2}{r}
\]
\end{tcolorbox}

\begin{tcolorbox}[title=สรุปบท: การหมุนของวัตถุแข็งเกร็ง ต่อ]
โมเมนต์ความเฉื่อย

\[
I=\sum_i m_i r_i^2
\]

พลังงานจลน์ของการหมุน

\[
K_{\text{rot}}=\frac{1}{2}I\omega^2
\]

การกลิ้งโดยไม่ไถล

\[
v_{\text{cm}}=\omega R
\]

\[
a_{\text{cm}}=\alpha R
\]

\[
K=\frac{1}{2}Mv_{\text{cm}}^2+\frac{1}{2}I_{\text{cm}}\omega^2
\]

ทอร์กและกฎการหมุน

\[
\tau=rF\sin\phi=Fd
\]

\[
\sum\tau=I\alpha
\]

สมดุลสถิต

\[
\sum F_x=0
\]

\[
\sum F_y=0
\]

\[
\sum\tau=0
\]
\end{tcolorbox}

% =========================================================
\section{ข้อสอบจริง สอวน. ที่ใช้ความรู้บทการหมุน}

\begin{tcolorbox}[title=คำแนะนำการใส่ภาพข้อสอบ]
ให้ตัดภาพข้อสอบจริงแล้วบันทึกตามชื่อไฟล์ที่กำหนดไว้ในแต่ละข้อ ถ้ายังไม่มีไฟล์ ภาพจะไม่แสดง แต่เอกสารยังคอมไพล์ได้
\end{tcolorbox}

% =========================================================
\subsection{ปี 2567 ข้อ 4: ล้อกลิ้งโดยไม่ไถลลงพื้นเอียง}

\vspace{1em}

\IfFileExists{img/posn67q4.png}{
\begin{center}
    \includegraphics[width=1\linewidth]{img/posn67q4.png}
\end{center}
}{
\begin{tcolorbox}[title=ตำแหน่งภาพที่แนะนำ]
ให้ตัดภาพข้อสอบปี 2567 ข้อ 4 แล้วบันทึกเป็น

\[
\texttt{img/posn67q4.png}
\]
\end{tcolorbox}
}

\begin{examplebox}{เฉลยละเอียด}
ล้อมีรัศมี

\[
R
\]

มวล

\[
M
\]

และโมเมนต์ความเฉื่อยรอบแกนผ่านศูนย์กลางมวล

\[
I
\]

กำลังกลิ้งโดยไม่ไถลลงพื้นเอียงมุม $\theta$ ต้องการหาขนาดของแรงเสียดทานที่พื้นเอียงกระทำต่อล้อ

เลือกทิศลงตามพื้นเอียงเป็นบวก

แรงตามแนวพื้นเอียงคือ

\[
Mg\sin\theta
\]

ลงตามพื้น และแรงเสียดทาน $f$ ขึ้นตามพื้น

ดังนั้นสมการการเลื่อนที่คือ

\[
Mg\sin\theta-f=Ma
\]

แรงเสียดทานเป็นแรงที่ทำให้ล้อหมุนรอบศูนย์กลางมวล ดังนั้น

\[
fR=I\alpha
\]

เพราะกลิ้งโดยไม่ไถล

\[
a=\alpha R
\]

จึงมี

\[
\alpha=\frac{a}{R}
\]

แทนในสมการทอร์ก

\[
fR=I\frac{a}{R}
\]

\[
f=\frac{I}{R^2}a
\]

แทนในสมการการเลื่อนที่

\[
Mg\sin\theta-\frac{I}{R^2}a=Ma
\]

\[
Mg\sin\theta=a\left(M+\frac{I}{R^2}\right)
\]

ดังนั้น

\[
a=\frac{Mg\sin\theta}{M+\frac{I}{R^2}}
\]

แรงเสียดทานคือ

\[
f=\frac{I}{R^2}a
\]

\[
f=
\frac{I}{R^2}
\left(
\frac{Mg\sin\theta}{M+\frac{I}{R^2}}
\right)
\]

จัดรูป

\[
\boxed{
f=\frac{I}{MR^2+I}Mg\sin\theta
}
\]

แรงเสียดทานมีทิศขึ้นตามพื้นเอียง แต่โจทย์ถามขนาด จึงตอบเฉพาะค่าบวก
\end{examplebox}

% =========================================================
\subsection{ปี 2561 ข้อ 10: ประตูกับแรงที่บานพับ}

\vspace{1em}

\IfFileExists{img/posn61q10.png}{
\begin{center}
    \includegraphics[width=1\linewidth]{img/posn61q10.png}
\end{center}
}{
\begin{tcolorbox}[title=ตำแหน่งภาพที่แนะนำ]
ให้ตัดภาพข้อสอบปี 2561 ข้อ 10 แล้วบันทึกเป็น

\[
\texttt{img/posn61q10.png}
\]
\end{tcolorbox}
}

\begin{examplebox}{เฉลยละเอียด}
ประตูหนัก

\[
W=90\,\si{N}
\]

ประตูกว้าง

\[
b=1.0\,\si{m}
\]

บานพับบนและล่างห่างกัน

\[
h=1.5\,\si{m}
\]

แรงที่บานพับบนกระทำต่อประตูมีขนาด

\[
F_{\text{บน}}=50\,\si{N}
\]

ต้องการหาแรงที่บานพับล่างกระทำต่อประตู

ให้แรงที่บานพับบนมีองค์ประกอบแนวระดับและแนวดิ่งเป็น

\[
H_{\text{บน}},\quad V_{\text{บน}}
\]

และแรงที่บานพับล่างเป็น

\[
H_{\text{ล่าง}},\quad V_{\text{ล่าง}}
\]

พิจารณาทอร์กรอบบานพับล่าง

น้ำหนักประตูกระทำที่กึ่งกลางความกว้าง ดังนั้นแขนโมเมนต์ของน้ำหนักคือ

\[
\frac{b}{2}=0.50\,\si{m}
\]

องค์ประกอบแนวระดับของแรงที่บานพับบนมีแขนโมเมนต์

\[
h=1.5\,\si{m}
\]

สมดุลทอร์กให้

\[
H_{\text{บน}}h=W\left(\frac{b}{2}\right)
\]

แทนค่า

\[
H_{\text{บน}}(1.5)=90(0.50)
\]

\[
H_{\text{บน}}=\frac{45}{1.5}
\]

\[
H_{\text{บน}}=30\,\si{N}
\]

เพราะขนาดแรงที่บานพับบนคือ $50\,\si{N}$

\[
F_{\text{บน}}^2=H_{\text{บน}}^2+V_{\text{บน}}^2
\]

\[
50^2=30^2+V_{\text{บน}}^2
\]

\[
V_{\text{บน}}=40\,\si{N}
\]

สมดุลแรงแนวดิ่ง

\[
V_{\text{บน}}+V_{\text{ล่าง}}=W
\]

\[
40+V_{\text{ล่าง}}=90
\]

\[
V_{\text{ล่าง}}=50\,\si{N}
\]

สมดุลแรงแนวระดับ

\[
H_{\text{บน}}+H_{\text{ล่าง}}=0
\]

ดังนั้นขนาดขององค์ประกอบแนวระดับที่บานพับล่างคือ

\[
|H_{\text{ล่าง}}|=30\,\si{N}
\]

ขนาดแรงลัพธ์ที่บานพับล่างคือ

\[
F_{\text{ล่าง}}
=
\sqrt{H_{\text{ล่าง}}^2+V_{\text{ล่าง}}^2}
\]

\[
F_{\text{ล่าง}}
=
\sqrt{30^2+50^2}
\]

\[
F_{\text{ล่าง}}
=
\sqrt{3400}
\]

\[
F_{\text{ล่าง}}
=
10\sqrt{34}\,\si{N}
\]

ประมาณเป็น

\[
\boxed{F_{\text{ล่าง}}\approx 58\,\si{N}}
\]
\end{examplebox}

% =========================================================
\subsection{ปี 2562 ข้อ 19: แท่งล้มบนพื้นลื่น}

\vspace{1.5em}

\IfFileExists{img/posn62q19.png}{
\begin{center}
    \includegraphics[width=1\linewidth]{img/posn62q19.png}
\end{center}
}{
\begin{tcolorbox}[title=ตำแหน่งภาพที่แนะนำ]
ให้ตัดภาพข้อสอบปี 2562 ข้อ 19 แล้วบันทึกเป็น

\[
\texttt{img/posn62q19.png}
\]
\end{tcolorbox}
}

\vspace{1em}

\begin{examplebox}{เฉลยละเอียด}
โจทย์นี้เป็นโจทย์ผสมระหว่างพลังงานกับการหมุนของวัตถุแข็งเกร็ง

แท่งโลหะผอมสม่ำเสมอยาว

\[
2l
\]

เริ่มเกือบอยู่ในแนวดิ่งบนพื้นลื่น และล้มลงจนจุดศูนย์กลางมวลชนพื้น

เพราะพื้นลื่น ไม่มีแรงเสียดทานจากพื้น ดังนั้นแรงภายนอกในแนวระดับเป็นศูนย์ จุดศูนย์กลางมวลจึงไม่มีความเร่งในแนวระดับ

ตอนเริ่มต้น ศูนย์กลางมวลอยู่สูงจากพื้น

\[
l
\]

ตอนศูนย์กลางมวลชนพื้น ความสูงของศูนย์กลางมวลเป็นศูนย์

ดังนั้นพลังงานศักย์โน้มถ่วงที่ลดลงคือ

\[
\Delta U=Mgl
\]

พลังงานนี้เปลี่ยนเป็นพลังงานจลน์ของการเลื่อนที่ของศูนย์กลางมวลและพลังงานจลน์ของการหมุน

\[
Mgl=\frac{1}{2}Mv_{\text{cm}}^2+\frac{1}{2}I_{\text{cm}}\omega^2
\]

สำหรับแท่งยาว \(2l\) หมุนรอบแกนผ่านศูนย์กลางมวลและตั้งฉากกับแท่ง

\[
I_{\text{cm}}=\frac{1}{12}M(2l)^2
\]

ดังนั้น

\[
I_{\text{cm}}=\frac{1}{3}Ml^2
\]

ขณะที่ศูนย์กลางมวลกำลังชนพื้น แท่งอยู่ในแนวราบ และจากเงื่อนไขเรขาคณิตของการหมุน จะได้

\[
v_{\text{cm}}=l\omega
\]

แทนลงในสมการพลังงาน

\[
Mgl=\frac{1}{2}M(l\omega)^2+\frac{1}{2}\left(\frac{1}{3}Ml^2\right)\omega^2
\]

\[
Mgl=\frac{1}{2}Ml^2\omega^2+\frac{1}{6}Ml^2\omega^2
\]

\[
Mgl=\frac{2}{3}Ml^2\omega^2
\]

ตัด \(M\)

\[
gl=\frac{2}{3}l^2\omega^2
\]

\[
l^2\omega^2=\frac{3gl}{2}
\]

แต่

\[
v_{\text{cm}}=l\omega
\]

ดังนั้น

\[
v_{\text{cm}}^2=\frac{3gl}{2}
\]

\[
\boxed{v_{\text{cm}}=\sqrt{\frac{3gl}{2}}}
\]
\end{examplebox}

% =========================================================
\subsection{ปี 2562 ข้อ 22: ล้อถูกดึงด้วยเชือกพันรอบเพลา}

\vspace{1.5em}

\IfFileExists{img/posn62q22.png}{
\begin{center}
    \includegraphics[width=1\linewidth]{img/posn62q22.png}
\end{center}
}{
\begin{tcolorbox}[title=ตำแหน่งภาพที่แนะนำ]
ให้ตัดภาพข้อสอบปี 2562 ข้อ 22 แล้วบันทึกเป็น

\[
\texttt{img/posn62q22.png}
\]
\end{tcolorbox}
}

\vspace{1em}

\begin{examplebox}{เฉลยละเอียด}
โจทย์นี้ใช้ความรู้ร่วมกันหลายบท คือ แรงลัพธ์ ทอร์ก และเงื่อนไขการกลิ้งโดยไม่ไถล

ให้ล้อมีมวล \(M\) รัศมีนอก

\[
R_2
\]

มีเพลารัศมี

\[
R_1
\]

และมีโมเมนต์ความเฉื่อย

\[
I
\]

แรง \(F\) ดึงเชือกที่ทำมุม \(\theta\) กับแนวระดับ

ให้ล้อกลิ้งไปทางขวา ดังนั้นต้องการให้ความเร่งของศูนย์กลางมวลเป็นบวก

\[
a>0
\]

เงื่อนไขกลิ้งโดยไม่ไถลคือ

\[
a=\alpha R_2
\]

หรือ

\[
\alpha=\frac{a}{R_2}
\]

สมการแรงตามแนวระดับคือ

\[
Ma=F\cos\theta-f
\]

สมการทอร์กรอบศูนย์กลางมวลคือ

\[
I\alpha=fR_2-FR_1
\]

แทน

\[
\alpha=\frac{a}{R_2}
\]

จะได้

\[
I\left(\frac{a}{R_2}\right)=fR_2-FR_1
\]

จัดรูปหา \(f\)

\[
fR_2=I\frac{a}{R_2}+FR_1
\]

\[
f=\frac{I}{R_2^2}a+\frac{FR_1}{R_2}
\]

นำไปแทนในสมการแรง

\[
Ma=F\cos\theta-\left(\frac{I}{R_2^2}a+\frac{FR_1}{R_2}\right)
\]

ย้ายพจน์ที่มี \(a\) ไปข้างเดียวกัน

\[
Ma+\frac{I}{R_2^2}a
=
F\cos\theta-\frac{FR_1}{R_2}
\]

\[
a\left(M+\frac{I}{R_2^2}\right)
=
F\left(\cos\theta-\frac{R_1}{R_2}\right)
\]

เพราะ

\[
M+\frac{I}{R_2^2}>0
\]

และ

\[
F>0
\]

ถ้าต้องการให้ล้อกลิ้งไปทางขวา ต้องมี

\[
\cos\theta-\frac{R_1}{R_2}>0
\]

ดังนั้น

\[
\cos\theta>\frac{R_1}{R_2}
\]

จึงได้เงื่อนไขของมุม

\[
\boxed{\theta<\cos^{-1}\left(\frac{R_1}{R_2}\right)}
\]
\end{examplebox}

% =========================================================
\subsection{ปี 2562 ข้อ 23: ความเร็วของจุดสีบนขอบล้อที่กลิ้ง}

\vspace{1.5em}

\IfFileExists{img/posn62q23.png}{
\begin{center}
    \includegraphics[width=1\linewidth]{img/posn62q23.png}
\end{center}
}{
\begin{tcolorbox}[title=ตำแหน่งภาพที่แนะนำ]
ให้ตัดภาพข้อสอบปี 2562 ข้อ 23 แล้วบันทึกเป็น

\[
\texttt{img/posn62q23.png}
\]
\end{tcolorbox}
}

\vspace{1em}

\begin{examplebox}{เฉลยละเอียด}
โจทย์นี้ใช้แนวคิดสำคัญของการกลิ้งโดยไม่ไถล คือความเร็วของจุดบนล้อเท่ากับผลรวมเวกเตอร์ของ

\[
\text{ความเร็วของศูนย์กลางมวล}
\]

และ

\[
\text{ความเร็วจากการหมุนรอบศูนย์กลางมวล}
\]

ล้อรัศมี \(R\) กลิ้งไปทางขวาด้วยความเร็วเชิงมุม

\[
\omega
\]

เงื่อนไขกลิ้งโดยไม่ไถลให้

\[
v_{\text{cm}}=R\omega
\]

ดังนั้น

\[
\vec{v}_{\text{cm}}=R\omega\hat{i}
\]

ให้จุดสีอยู่บนขอบล้อ โดยรัศมีที่ลากจากศูนย์กลางไปยังจุดสีทำมุม \(\theta\) กับแนวดิ่งขึ้น ตามรูป

ความเร็วจากการหมุนของจุดสีเทียบกับศูนย์กลางมีองค์ประกอบ

\[
\vec{v}_{\text{rot}}
=
R\omega\cos\theta\hat{i}
-
R\omega\sin\theta\hat{j}
\]

ดังนั้นความเร็วของจุดสีเทียบกับพื้นคือ

\[
\vec{v}
=
\vec{v}_{\text{cm}}+\vec{v}_{\text{rot}}
\]

\[
\vec{v}
=
R\omega\hat{i}
+
R\omega\cos\theta\hat{i}
-
R\omega\sin\theta\hat{j}
\]

\[
\vec{v}
=
R\omega(1+\cos\theta)\hat{i}
-
R\omega\sin\theta\hat{j}
\]

ขนาดของความเร็วคือ

\[
v=
\sqrt{
\left[R\omega(1+\cos\theta)\right]^2
+
\left[-R\omega\sin\theta\right]^2
}
\]

ดึง \(R\omega\) ออก

\[
v=R\omega
\sqrt{(1+\cos\theta)^2+\sin^2\theta}
\]

ขยายกำลังสอง

\[
(1+\cos\theta)^2+\sin^2\theta
=
1+2\cos\theta+\cos^2\theta+\sin^2\theta
\]

เพราะ

\[
\sin^2\theta+\cos^2\theta=1
\]

จึงได้

\[
(1+\cos\theta)^2+\sin^2\theta
=
2+2\cos\theta
\]

ดังนั้น

\[
v=R\omega\sqrt{2+2\cos\theta}
\]

ใช้เอกลักษณ์

\[
2+2\cos\theta=4\cos^2\left(\frac{\theta}{2}\right)
\]

จึงได้

\[
\boxed{v=2R\omega\cos\left(\frac{\theta}{2}\right)}
\]
\end{examplebox}

% =========================================================
\subsection{ปี 2566 ข้อ 28: แผ่นสี่เหลี่ยมแขวนและมวลถ่วง}

\vspace{1.5em}

\IfFileExists{img/posn66q28.png}{
\begin{center}
    \includegraphics[width=1\linewidth]{img/posn66q28.png}
\end{center}
}{
\begin{tcolorbox}[title=ตำแหน่งภาพที่แนะนำ]
ให้ตัดภาพข้อสอบปี 2566 ข้อ 28 แล้วบันทึกเป็น

\[
\texttt{img/posn66q28.png}
\]
\end{tcolorbox}
}

\vspace{1em}

\begin{examplebox}{เฉลยละเอียด}
โจทย์นี้เป็นสมดุลสถิตของวัตถุแข็งเกร็ง ใช้แนวคิดว่า ถ้าวัตถุถูกแขวนที่จุด \(A\) แล้วหยุดนิ่ง จุดศูนย์กลางมวลรวมของระบบต้องอยู่ในแนวดิ่งใต้จุดแขวน

แผ่นสี่เหลี่ยมมุมฉากมีมวล

\[
M
\]

ยาว

\[
L
\]

กว้าง

\[
b
\]

จุด \(A\) เป็นจุดกึ่งกลางของขอบบน และมีมวล \(m\) แขวนที่มุมล่าง \(B\) ตามรูป

เลือกแกนที่ติดกับแผ่น โดยให้แกนตามแนวยาวของแผ่นเป็นแกน \(x'\) และแกนลงตามความกว้างของแผ่นเป็นแกน \(y'\)

เมื่อแผ่นเอียงทำมุม \(\theta\) กับแนวระดับ ตำแหน่งแนวระดับของศูนย์กลางมวลของแผ่นเทียบกับจุด \(A\) คือ

\[
x_M=\frac{b}{2}\sin\theta
\]

ตำแหน่งแนวระดับของมวล \(m\) ที่จุด \(B\) เทียบกับจุด \(A\) คือ

\[
x_m=b\sin\theta-\frac{L}{2}\cos\theta
\]

เงื่อนไขสมดุลคือ จุดศูนย์กลางมวลรวมต้องอยู่ตรงใต้จุดแขวน ดังนั้นโมเมนต์ของน้ำหนักรอบจุด \(A\) ต้องเป็นศูนย์

\[
Mgx_M+mgx_m=0
\]

ตัด \(g\)

\[
Mx_M+mx_m=0
\]

แทนค่า \(x_M\) และ \(x_m\)

\[
M\left(\frac{b}{2}\sin\theta\right)
+
m\left(b\sin\theta-\frac{L}{2}\cos\theta\right)
=0
\]

กระจายพจน์

\[
\frac{Mb}{2}\sin\theta
+
mb\sin\theta
-
\frac{mL}{2}\cos\theta
=0
\]

ย้ายพจน์

\[
\frac{Mb}{2}\sin\theta
+
mb\sin\theta
=
\frac{mL}{2}\cos\theta
\]

ดึงตัวร่วมฝั่งซ้าย

\[
\frac{b}{2}(M+2m)\sin\theta
=
\frac{mL}{2}\cos\theta
\]

คูณทั้งสองข้างด้วย \(2\)

\[
b(M+2m)\sin\theta
=
mL\cos\theta
\]

หารด้วย \(b(M+2m)\cos\theta\)

\[
\tan\theta
=
\frac{mL}{b(M+2m)}
\]

ดังนั้น

\[
\boxed{
\tan\theta=\frac{mL}{b(M+2m)}
}
\]
\end{examplebox}